\chapter{Resolution, useful accuracy and the cost of an experiment}
\label{mathematics-resolution}

Nimbus's aim includes high-fidelity atmospheric experiments on ordinary
workstations and MacBooks. Efficient algorithms and GPU acceleration can
change which experiments are practical on a given machine. The mathematical
question is what accuracy an experiment needs and which changes actually buy
that accuracy. A smaller grid spacing, a faster timestep implementation and a
more detailed rendered image are three different improvements. This chapter
connects the physical question to discretization, computational cost and the
evidence needed to increase resolution deliberately.

The discussion supplies planning models and experiment-design guidance. It
does not assign a universal maximum grid to a laptop, recommend one spacing
for all storms, or claim that Nimbus is already fully optimized. The
\href{ARCHITECTURE.pdf}{architecture guide} explains how to measure memory,
data movement and execution, and what distributed GPU execution would require.

\section{Define fidelity through a question and an observable}

Let $X(t)$ denote the complete model state and $J[X]$ a quantity of interest.
Examples include domain-integrated water, accumulated precipitation over a
specified area and interval, an updraft-volume distribution, a cold-pool
propagation speed, or a spectrum of resolved kinetic energy. Define the region,
time interval, units and averaging operation before choosing the grid. A
single maximum of vertical vorticity and an area-averaged circulation are not
equivalent targets and need not converge at the same rate.

For example, a dry-mass-weighted mean mixing ratio over a fixed physical region
$\mathcal R$ is
\[
 J_q(t)=\frac{\displaystyle\sum_{C\subset\mathcal R}M_C(t)q_C(t)}
              {\displaystyle\sum_{C\subset\mathcal R}M_C(t)},
 \qquad M_C=\rho_{d,C}V_C.
\]
This expression assumes the region is a union of owned physical cells.
Partially intersected cells need an explicit overlap rule. Terrain volumes,
moving domains and coarse/fine coverage change those weights; storage halos
never contribute merely because they contain finite values. Comparing runs
on different grids requires a common physical region rather than the same
integer array slice.

High fidelity therefore has several components: the appropriate governing
approximations, credible initial and boundary conditions, suitable closure
assumptions, sufficient spatial and temporal resolution, and a numerical
implementation that preserves the declared contracts. The objective is to
make that combination affordable locally. Optimization can remove waste
without changing the question; changing the question to fit the machine must
be a deliberate scientific choice.

\section{Grid spacing is not the smallest trustworthy feature}

Sampling a wavelength at two points is a representability limit, not a promise
that an evolution scheme propagates that wavelength accurately. To see the
distinction, consider the teaching example of a centered derivative on a
uniform collocated grid,
\[
 D_h f_j=\frac{f_{j+1}-f_{j-1}}{2h},\qquad f_j=e^{ikjh}.
\]
Substitution gives
\[
 D_h f_j=i\frac{\sin(kh)}{h}f_j,
 \qquad \frac{k_{\mathrm{effective}}}{k}
       =\frac{\sin(kh)}{kh}.
\]
At wavelengths $\lambda=4h$ and $8h$, this ratio is approximately 0.637 and
0.900. At the two-cell alternating mode it is zero. The formula illustrates
why a visibly sampled pattern can have a badly represented derivative. It is
not the transfer function of Nimbus's full scheme: face staggering,
reconstruction, upwinding, limiters, dissipation and time integration change
the result. In particular, a face-centered gradient has a different symbol
from this collocated example.

An effective-resolution study must examine the combined algorithm. Wave
propagation and decay tests expose phase and amplitude errors; spectra and
structure statistics expose where an evolving turbulent solution loses useful
scale content. Skamarock's
\href{https://doi.org/10.1175/MWR2830.1}{kinetic-energy-spectrum study}
is relevant methodological context. It does not establish Nimbus's effective
resolution or justify transferring another model's cells-per-wavelength
estimate without testing the Nimbus configuration.

The same issue matters in a vorticity rendering. Curl uses spatial derivatives
of velocity, and its apparent small structures can be sensitive to filtering,
stagger-to-center reconstruction and the derivative operator. A smooth GPU
image can make a field easier to inspect while leaving its physical resolving
power unchanged. Report the velocity source, curl definition, units,
sampling, thresholds and display transformation with such a figure. The
presence of fine visible structure is useful evidence to investigate, but its
resolved scales must be established from the underlying fields and experiment.

\section{Refinement changes the resolved and unresolved partition}

In a filtered-flow interpretation, a representative subgrid stress is
\[
 \tau_{ij}=\overline{u_i u_j}-\overline{u}_i\,\overline{u}_j.
\]
This expression states that filtering a product differs from multiplying
filtered velocities. It is an illustrative incompressible form; a compressible
closure may use density-weighted filtering and Nimbus's selected scheme has
its own precise conventions. Refining the grid changes what is explicitly
represented and what the closure must supply. It does not simply produce more
samples of an otherwise unchanged turbulent model.

For an eddy-viscosity example, $K\sim(C\Delta)^2|S|$. Both the filter scale
$\Delta$ and resolved strain $|S|$ can change on refinement. Holding the
coefficient $C$ fixed therefore does not hold the realized mixing fixed.
The relationship between vertical and horizontal resolution also matters:
a thin near-surface layer and a wide horizontal cell do not have the same
numerical or closure behavior as an isotropic cell of equal volume.

Bryan, Wyngaard and Fritsch's
\href{https://journals.ametsoc.org/view/journals/mwre/131/10/1520-0493_2003_131_2394_rrftso_2.0.co_2.xml}{study
of deep moist convection} motivates examining the scale assumptions of an
LES closure alongside resolution. Its results are not a universal grid cutoff
for every Nimbus experiment. Bulk microphysics, cloud optical properties,
surface exchange and initial-condition uncertainty remain modeled even after
important turbulent motions become resolved.

Refinement also needs compatible inputs. Preserve the physical size and
amplitude of an initial perturbation rather than its radius in cells. Sample
the same sounding and terrain with documented reconstruction. Keep sponge
thickness and forcing regions in physical units where that defines the
experiment. Recheck the realized base-state balance. Otherwise a purported
resolution comparison changes the triggering or boundaries at the same time.

\section{Why halving spacing can cost much more than eight times}

For a fixed rectangular physical domain, ignoring halos and staggered end
planes, the cell count is approximately
\[
 N\simeq\frac{L_xL_yL_z}{\Delta x\,\Delta y\,\Delta z}.
\]
Refining each direction by factors $(r_x,r_y,r_z)$ increases the cell count by
$r_xr_yr_z$. The time cost also depends on the timestep. A schematic explicit
wave/advection restriction is
\[
 \Delta t\left(
 \frac{|u|+c_x}{\Delta x}
 +\frac{|v|+c_y}{\Delta y}
 +\frac{|w|+c_z}{\Delta z}\right)\lesssim C_{\mathrm{scheme}},
\]
where the relevant characteristic speeds and stability constant depend on the
operator and integrator. This is a dimensional planning relation, not Nimbus's
exact timestep controller. Moving-grid transport uses velocity relative to
the grid, and mapped coordinates require the corresponding metric factors.
Split-explicit integration separates slow-step and
acoustic restrictions; the cost can increase through more outer steps, more
fast substeps, or both. Terrain metrics and local thin layers can determine the
restrictive physical spacing.

If all spacings shrink by $r$, characteristic speeds stay comparable, the
explicit timestep scales as $1/r$, and work per cell per step remains comparable,
the same physical duration requires approximately $r^4$ work in three
dimensions. Doubling linear resolution gives eight times the cells and about
sixteen times the work under those assumptions. Horizontal-only refinement
gives roughly $r^3$ work if the horizontal CFL restriction controls the step;
if an unchanged vertical restriction already dominates, the initial increase
can instead be closer to $r^2$. These are distinct cases, not conflicting rules.

For an explicit constant-diffusivity operator, a representative restriction is
\[
 K\Delta t\left(\Delta x^{-2}+\Delta y^{-2}+\Delta z^{-2}\right)
 \lesssim C_{\mathrm{diff}}.
\]
If this becomes controlling and $K$ stays fixed, isotropic refinement can make
the work scale as $r^5$. A grid-dependent SGS coefficient, implicit solve or
subcycling changes that estimate. Positivity restrictions, sedimentation and
microphysics can add costs that do not obey the simple wave model. For example,
increasing the number of cloudy cells can change the work distribution even
when the total cell count scales regularly.

\begin{longtable}{@{}>{\raggedright\arraybackslash}p{0.36\textwidth}>{\centering\arraybackslash}p{0.18\textwidth}>{\raggedright\arraybackslash}p{0.37\textwidth}@{}}
\toprule
Illustrative change & State-size factor & Work factor for fixed duration \\
\midrule
\endhead
Halve all three spacings & $8$ & About $16$ if wave CFL controls \\
Halve horizontal spacings only & $4$ & About $8$ if horizontal CFL controls \\
Halve all three spacings, fixed explicit diffusivity controls & $8$ & About $32$ \\
Double physical domain lengths at unchanged spacing & $8$ & About $8$ if step
size and per-cell work stay comparable \\
\bottomrule
\end{longtable}

The factors exclude changes in communication, cache behavior, output,
initialization and iteration counts. They explain why optimization has large
scientific value: removing repeated scans, avoidable synchronization or
unneeded copies can make the next refinement practical without changing its
equations. They also explain why a speedup on one grid cannot be extrapolated
unchanged across a memory-capacity boundary. Actual complete-run measurements
must determine the local hardware's practical range.

The same model quantifies an optimization opportunity. Suppose a complete
matched workload becomes $S$ times faster, and its refinement cost remains
approximately proportional to $h^{-4}$. At a fixed time budget, that gain could
support a linear refinement factor $r\simeq S^{1/4}$, provided memory and other
constraints still fit. An illustrative sixteen-fold gain corresponds to
halving all three spacings for about the original elapsed time. This is not
a measured Nimbus speedup. The refined state still has eight times as many
cells, so reducing unnecessary live storage can be as consequential as faster
arithmetic. Alternatively, the same gain can support more realizations at an
already adequate resolution. This is why optimization changes the accessible
science on ordinary hardware rather than merely changing a benchmark number.

\section{Memory and output belong in the same resolution calculation}

With $F$ simultaneously live scalar-equivalents and $b$ bytes per scalar, a
first payload estimate is $bFN$. It omits unequal field shapes, geometry,
halos, stage duplication, private scheme state and output snapshots unless
those have explicitly been included in $F$. Storage precision also differs
from arithmetic precision. A float32 field with double-precision temporary
work cannot be budgeted as though every live representation uses four bytes.

For a halo depth $h$ around an owned local tile, the allocation ratio is
\[
 R_{\mathrm{storage}}=
 \frac{(n_x+2h)(n_y+2h)(n_z+2h)}{n_xn_yn_z}.
\]
This simple rectangular example shows why small distributed tiles can have
high overhead. Physical boundary layouts and C-grid fields change the exact
count. Halos increase storage and communication, but not the physical
atmospheric inventory used in a conservation check.

If $F_o$ fields are saved for physical duration $T$ at interval $\tau_o$,
an uncompressed payload estimate is
\[
 B_{\mathrm{history}}\simeq
 bF_oN\left(1+\left\lfloor T/\tau_o\right\rfloor\right).
\]
It assumes an initial frame and fixed-size full fields, and excludes metadata,
checkpoints and predictive dependencies. At fixed physical output cadence,
isotropic refinement multiplies this payload by $r^3$. Saving every fixed
number of shrinking timesteps can add another factor of $r$. A comparison
should preserve the intended physical sampling rather than accidentally
turning a resolution experiment into a different I/O workload.

Compression changes storage and transfer costs according to its actual mode
and data. It does not make the uncompressed computational working set vanish.
Some read paths can decode a bounded window; others still need a larger source
working set. A measured compression ratio for one smooth field is not an
admission guarantee for all fields or storm stages.

\section{A refinement study compares the same mathematical experiment}

For a smooth, sufficiently resolved verification problem, suppose
$J(h)=J_*+Ch^p+o(h^p)$ with a stable leading error coefficient. For three
grids $h,h/r,h/r^2$ with $r>1$, an observed order is
\[
 p_{\mathrm{obs}}=
 \frac{\log\left|\dfrac{J(h)-J(h/r)}{J(h/r)-J(h/r^2)}\right|}{\log r}.
\]
If the differences are meaningful and the asymptotic model applies, the
corresponding extrapolation is
\[
 J_*\approx J(h/r^2)+
 \frac{J(h/r^2)-J(h/r)}{r^{p_{\mathrm{obs}}}-1}.
\]
These formulas are conditional estimates. Near-zero differences, nonmonotone
behavior, phase shifts or changing physical closure can invalidate the
interpretation. NASA's
\href{https://www.grc.nasa.gov/www/wind/valid/tutorial/spatconv.html}{spatial
convergence tutorial} provides verification context. A formally high-order
reconstruction does not establish the order of a complete moist, limited,
split-coupled simulation.

Start with manufactured or analytical problems that isolate spatial and
temporal errors. Refine time at a fixed grid, and refine space with temporal
error controlled. Test hydrostatic balance, wave propagation, translated
tracers, physical inventories and the selected source laws. For common-grid
comparisons, restrict conserved inventories with their physical weights;
arithmetic interpolation of mixing ratios is not generally conservative.
Retain original fine-grid peaks as separate quantities if they are of interest.

Long nonlinear convective runs need additional interpretation. Small admissible
changes can shift storm timing or location and eventually produce different
realizations. Pointwise agreement at a late time is then a different question
from agreement of precipitation, transport, spectra or event distributions.
Compare those statistics on declared common scales and intervals. If a
storm-relative alignment is used, report it alongside the unaligned comparison;
it must not hide a propagation error. Neither visual similarity nor divergence
alone settles numerical correctness.

\section{Choose ensemble size and resolution together}

Suppose an experiment estimates a mean diagnostic from $m$ independent,
identically distributed realizations with finite variance $\sigma^2$. Its
sampling standard error scales as $\sigma/\sqrt m$. This does not describe
discretization bias, model-form error or a single deterministic simulation.
Snapshots from one trajectory are usually correlated and must not be counted
as independent realizations simply because many frames were saved.

A fixed compute budget can therefore be spent on finer individual runs, more
realizations, longer sampling or a larger physical domain. Which choice helps
depends on the quantity of interest and its dominant error. Reproducible
perturbation seeds help paired comparisons, but do not create independence
or validate the perturbation model. The local-hardware objective includes
making ensembles practical, not only producing the finest isolated frame.
Independent realizations are also a natural unit for spreading work across
additional computers without domain-exchange traffic.

\section{Keep numerical, physical and representation errors distinguishable}

An experiment record should separate spatial and temporal discretization,
closure and coefficient assumptions, initial and boundary uncertainty,
sampling variability, floating-point effects, and saved-data transformations.
Their effects interact; the list is not a theorem that their measured errors
can always be added as independent scalar terms. It identifies what to vary
and what to hold fixed in a controlled comparison.

For example, if saved velocity values have a pointwise error bound
$|e_j|\leq\epsilon$, a centered derivative reconstructed from those saved
values has representation error bounded by
\[
 \left|\frac{e_{j+1}-e_{j-1}}{2h}\right|\leq\frac{\epsilon}{h}.
\]
A finer grid can amplify a fixed saved-field error in a derivative diagnostic.
This bound assumes the pointwise premise and this specific stencil; it does
not assert that every compression mode supplies that premise. Vorticity and
other derivative products need their own error assessment. Lossy visualization
output is also different from lossy prognostic state or restart, which can
affect subsequent evolution.

The useful result of this process is a declared experiment envelope: the
observable, physical setup, supported configuration, grid and timestep,
input identities, output transformations, measured resource cost and observed
sensitivity. That envelope can demonstrate demanding local experiments where
the evidence supports them, and identify when additional optimization, memory,
ensemble throughput or distributed execution is needed. Hardware placement
follows the experiment and measurements; it is not a substitute for them.
